\section{需求分析}

本节全面分析新一代全尺寸金融市场仿真器的需求，包括目标用户、功能性需求、非功能性需求以及约束条件。通过明确这些需求，可为后续系统设计和模块划分提供依据。

\subsection{用户场景与利益相关者}

系统的主要用户包括以下类型：
\begin{itemize}
  \item \textbf{科研人员：}需要一个可控、可重复的实验环境，用于研究市场微观结构、市场行为、市场政策及其影响，进行回溯实验和对比分析。
  \item \textbf{量化交易策略开发者：}希望在无需真实交易的情况下测试和验证新的交易策略或算法，包括强化学习代理；需要实时插入策略和监控效果的能力。
  \item \textbf{监管机构与政策分析师：}用于模拟政策干预（如交易税、熔断机制、涨跌停制度）对市场流动性和波动性的影响，从而评估政策的有效性。
  \item \textbf{教育者与学生：}可用于教学目的，帮助理解订单簿机制、市场微观结构和智能体交互。
\end{itemize}

\subsection{功能性需求}

\subsubsection{多资产与跨市场支持}

\begin{itemize}
  \item 仿真器应支持同时模拟多只股票或其他资产，允许代理在多个市场或多个交易品种之间下达订单。每个资产需要独立维护订单簿，但也要能反映市场间的联动和流动性传导现象。
  \item 交易所模块应支持多市场配置，能够灵活定义交易时间、最小价格变动单位、撮合规则等，并允许各市场之间的信息互通（例如跨市场套利）。
\end{itemize}

\subsubsection{多时段及跨日机制}

\begin{itemize}
  \item 系统应不仅支持日内连续双向拍卖，还要提供集合竞价模块，用于模拟开盘、收盘或中场集合竞价，并允许用户灵活配置。
  \item 仿真应支持跨日状态保持，包括隔夜持仓、订单续存以及跨日价格调整，为研究跨日策略和交易时间切换提供支持。
\end{itemize}

\subsubsection{订单簿管理与撮合}

\begin{itemize}
  \item 订单簿模块应实现高效的数据结构，支持插入、删除、匹配和查询等操作，能够在高频率下维持良好性能。
  \item 匹配逻辑需严格遵循价格—时间优先原则，支持限价单、市场单、撤单等多种订单类型，并提供部分成交、滑点处理等机制。
\end{itemize}

\subsubsection{代理模型与自校准机制}

\begin{itemize}
  \item 系统应提供丰富的基础代理库，包括市场造市商、被动投资者、零智商代理、启发式学习代理、强化学习代理等；用户能够继承并扩展这些类实现自定义策略。
  \item 引入内部自校准机制，使背景代理能够根据历史行情或实时统计特征自动调整行为，实现与真实市场统计属性（例如波动率、交易量分布）的动态对齐，减少外部迭代调参的成本。
  \item 代理应能够感知并响应市场动态，包括多资产信息、交易费用变化、政策干预等，并具备自适应学习能力。
\end{itemize}

\subsubsection{数据管理与接口}

\begin{itemize}
  \item 系统应支持从多源数据加载历史行情，包括交易记录、行情快照、企业事件等，并提供统一的数据访问接口供代理查询。
  \item 支持生成合成数据，如基于统计模型或随机过程的行情序列，以便模拟极端情形或未曾发生过的场景。
  \item 提供灵活的配置文件或 API，用于指定实验参数、代理数量、市场规则、随机种子等。
\end{itemize}

\subsubsection{交互与可视化}

\begin{itemize}
  \item 支持仿真过程中实时注入新代理、调整参数或实施政策变化，并及时反馈结果，适用于“what-if”分析。
  \item 提供可视化界面或日志工具，显示订单簿深度、成交记录、价格曲线、代理持仓和盈亏等信息，以便用户监控仿真过程。
  \item 支持事件订阅和回调机制，使外部模块（如强化学习算法）能在仿真运行中接收市场状态、发出决策并得到执行结果。
\end{itemize}

\subsection{非功能性需求}

\subsubsection{性能与可扩展性}

\begin{itemize}
  \item 目标系统需要在高并发场景下保持高吞吐量，能够处理全市场级别的数据规模。例如，单一市场在高峰期可能出现超过 $300{,}000$ 笔订单/秒，需要仿真器的订单处理吞吐量不少于该数量级。
  \item 应支持在多核 CPU 或 GPU 环境下并行运行，提升多资产与多智能体仿真能力，并控制内存占用。
  \item 仿真结果应具有高度一致性和可重复性，在给定相同随机种子和配置的情况下，能够生成完全相同的输出，用于公平比较不同策略。
\end{itemize}

\subsubsection{可靠性与准确性}

\begin{itemize}
  \item 系统应严格遵循现实市场的交易规则（例如价格单位、时间优先规则、不同订单类型），以确保模拟结果具有可信度。
  \item 自校准机制应避免过度拟合，既要逼近历史统计特征，又要保留一定随机性，保证实验的可探索性。
\end{itemize}

\subsubsection{可维护性与可扩展性}

\begin{itemize}
  \item 软件架构应模块化，具备清晰的接口定义和依赖关系，便于后续在不同模块（如代理、撮合引擎、数据接口）上独立扩展或替换实现。
  \item 提供完善的文档，包括模块说明、使用方法、配置示例和 API 说明，帮助用户快速上手并进行二次开发。
\end{itemize}

\subsubsection{可用性}

\begin{itemize}
  \item 提供人性化的配置方式，支持通过 YAML、JSON 或 Python 脚本进行实验配置，同时允许命令行参数覆盖部分配置，方便批量实验。
  \item 支持在不同操作系统环境（如 Linux、MacOS、Windows）下部署和运行，降低使用门槛。
\end{itemize}

\subsection{约束与假设}

\begin{itemize}
  \item \textbf{真实市场约束：} 仿真器需要符合真实市场的监管法规、交易时间、价格最小变动单位等；某些特殊交易机制（如特殊的中止机制或暗池交易）可视需求后续扩展。
  \item \textbf{数据可用性：} 自校准和历史回放功能依赖高质量的历史交易数据和行情快照；对于无法获取的数据，需采用合成或估计方法。
  \item \textbf{资源限制：} 对于全市场仿真，需要大量计算资源支持高吞吐量和长时间仿真，故系统需支持在分布式或云环境中部署，用户需考虑资源配置。
  \item \textbf{用户技能假设：} 假设用户具备一定的编程基础和金融知识，能够理解基本的交易规则和代理模型；针对初学者需提供示例和指导。
\end{itemize}

综上所述，需求分析明确了新一代金融市场仿真器在功能、性能、可用性和扩展性方面需要满足的条件。这些需求将为后续的系统架构设计和模块划分提供指导。